% No novelty claim is made for the quotient or character-sum ingredients.
\subsection{Full-field norm quotients with almost every challenge}
\label{subsec:native-norm-quotients}

The quotient construction of Krachun, Kazanin, and Hab\"ock~\cite[Appendix~A]{kkh} can be applied
directly in the quotient variable.  A pole outside its image suffices; it need
not be the image of an evaluation-field point.  Combining this observation
with a character-sum estimate gives the following native-field specialization.

\begin{lemma}\label{lem:fixed-cardinality-character-products}
Let $\mathcal A$ be an $s$-element subset of a finite abelian group
$H$ of order $M$.  Suppose every nontrivial character $\chi$ of $H$
satisfies $|\sum_{a\in\mathcal A}\chi(a)|\le B<s$.
For $0<r<s$, put $\theta=r/s$ and
\[
 \varepsilon=(s+1)\exp\bigl(-\theta(1-\theta)(s-B)\bigr).
\]
If $(M-1)\varepsilon<1$, every element of $H$ is the product of
an $r$-element subset of $\mathcal A$.
\end{lemma}
\begin{proof}
Fix a nontrivial character and write $v_a=\chi(a)$.  For $t>0$ and
real $\phi$, the identity
\[
 |1+te^{i\phi}v_a|^2
 =(1+t)^2-2t\bigl(1-\operatorname{Re}(e^{i\phi}v_a)\bigr)
\]
and $1-u\le e^{-u}$ give
\[
 \left|\prod_{a\in\mathcal A}(1+te^{i\phi}v_a)\right|
 \le (1+t)^s\exp\left(-\frac{t(s-B)}{(1+t)^2}\right).
\]
Cauchy's coefficient bound, with $t=r/(s-r)$, bounds the absolute
value of the degree-$r$ coefficient by
$t^{-r}(1+t)^s\exp(-\theta(1-\theta)(s-B))$.
The value $r$ is a mode of the binomial distribution with parameters
$s,r/s$, so
$t^{-r}(1+t)^s\le(s+1)\binom sr$.
Consequently every nontrivial character contributes at most
$\varepsilon\binom sr$ in absolute value.  Character orthogonality
therefore bounds the number of $r$-subsets with any prescribed product
from below by
\[
 \frac1M\binom sr\bigl(1-(M-1)\varepsilon\bigr)>0.
\]
\end{proof}

\begin{theorem}\label{thm:native-norm-quotient}
Fix $0<\rho<1$ and an integer $d\ge2$.
For every sufficiently large prime $p$, put
\[
 q=p^d,\qquad E=\mathbb F_q,\qquad
 m=\frac{q-1}{p-1},\qquad J=\lfloor\rho q\rfloor.
\]
There are words $f,g:E\to E$ such that
\[
 \operatorname{agr}_J(f)=\operatorname{agr}_J(g)
 =\operatorname{CA}_J(f,g)=J+m-1,
\]
and every $\lambda\in E\setminus\{0\}$ satisfies
\[
 \operatorname{agr}_J(f+\lambda g)\ge J+2m-1.
\]
Here the code has strict degree bound $J$, the capacity margin is
$(2m-1)/q=\Theta(p^{-1})$, and the source-agreement gap is
$m/q=\Theta(p^{-1})$.

For an explicit sufficient condition, write
$J-1=(r-2)m+w$ with $0\le w<m$, and set $\theta=r/(p-2)$.
It is enough that $2\le r<p-2$ and
\[
 \log\bigl((p^d-2)(p-1)\bigr)
 <\theta(1-\theta)\bigl(p-4-(d-1)\sqrt p\bigr).
\]
\end{theorem}

\begin{proof}
Choose $b\in E$ of degree exactly $d$ over $\mathbb F_p$.
Write $N(X)=X^m$, so that $N$ is the
norm on $E$, with image $\mathbb F_p$.  Fix $a_0\in\mathbb F_p^*$, choose
a $w$-element subset $B\subset N^{-1}(a_0)$, and put
$R(X)=\prod_{x\in B}(X-x)$.  In particular $R$ is monic of degree $w<m$.
Define
\[
 f(X)=R(X)\frac{N(X)^r-b^r}{N(X)-b},\qquad
 g(X)=-\frac{R(X)}{N(X)-b}.
\]
The first expression is a monic polynomial of degree
$(r-1)m+w=J+m-1$; the denominator in the second expression never vanishes
on $E$.

For every polynomial $h$ of degree less than $J$, root counting bounds the
agreement of $f$ with $h$ by $J+m-1$.  The same bound applies to $g$:
its agreement equation is
$R+(N-b)h=0$, a nonzero polynomial of degree at most $J+m-1$.
It is nonzero because $\deg R<m$.  These two bounds are attained
simultaneously.  Choose any $r-1$ distinct tags in
$\mathbb F_p^*\setminus\{a_0\}$, interpolate both functions
$(Y^r-b^r)/(Y-b)$ and $-1/(Y-b)$ on those tags by polynomials of degree
at most $r-2$, substitute $Y=N(X)$, and multiply by $R(X)$.
Both resulting witnesses have degree at most $(r-2)m+w=J-1$.
They match on the selected $r-1$ norm fibers and on $B$, a disjoint set
of size $(r-1)m+w=J+m-1$.  This proves all three stated equalities.

It remains to represent every element of $E^*$ as a product
\[
 \prod_{a\in S}(b-a),\qquad
 S\subset\mathbb F_p^*\setminus\{a_0\},\quad |S|=r.
\]
For every nontrivial multiplicative character $\chi$ of $E^*$,
Katz's affine-line estimate~\cite[Theorem~1]{katz1989-character-sums}
bounds $|\sum_{a\in\mathbb F_p}\chi(b-a)|$ by $(d-1)\sqrt p$.
Deleting $0,a_0$ gives a bound $(d-1)\sqrt p+2$ on the remaining $p-2$ terms.
Lemma~\ref{lem:fixed-cardinality-character-products} consequently proves
the required surjectivity under the displayed sufficient condition.
For fixed $\rho,d$,
$\theta\to\rho$, so this condition holds for all sufficiently large $p$.

For $\lambda\ne0$, choose such an $S$ with product $-\lambda$, and set
\[
 V_S(Y)=\prod_{a\in S}(Y-a),\qquad
 P_S(Y)=Y^r-V_S(Y),\qquad
 h_S(X)=R(X)\frac{P_S(N(X))-P_S(b)}{N(X)-b}.
\]
The quotient is polynomial and
$\deg h_S\le w+(r-2)m=J-1$.  Since
$P_S(b)=b^r-V_S(b)=b^r+\lambda$, direct subtraction gives
\[
 f(X)+\lambda g(X)-h_S(X)
 =\frac{R(X)V_S(N(X))}{N(X)-b}.
\]
Its zeros are exactly $B$ and the $r$ norm fibers indexed by $S$.
These sets are disjoint and have total size
$w+rm=J+2m-1$.  In particular $X=0$ supplies no hidden extra agreement.
\end{proof}

In particular, put $T=J+2m-1$ and $\delta=1-T/q$.
Both sources have distance exactly $\delta+m/q$ from the code,
but every affine combination $(1-t)f+tg$ with $t\notin\{0,1\}$
has distance at most $\delta$. Indeed, divide by $1-t$ and apply
the theorem with $\lambda=t/(1-t)\ne0$.
Thus the nearby fraction on this affine line is exactly $1-2/q$.
The normalized common-agreement loss $m/q$ is asymptotically half
the capacity margin $(2m-1)/q$.

For $d=5$, this gives a larger margin, of order $q^{-1/5}$, than
the order-$q^{-2/5}$ margin of
Theorem~\ref{thm:growing-characteristic-native}, and both sources
are individually far. It does not preserve that theorem's exact
common agreement $J$: here common agreement is $J+m-1$.
Both margins are smaller than the inverse-logarithmic margin in~\cite{kkh}.
The conclusion is a native-field coverage refinement of standard
quotient and character-sum ingredients, not a better capacity-margin
benchmark. The alphabet is an extension field, $p\ll J$, and the
evaluation domain has all $p^d$ points. No construction on the short
domains used by the practical benchmark follows from this theorem.
